\documentclass[12pt]{article}
\usepackage{amsmath, amssymb}
\usepackage[margin=1in]{geometry}
\setlength{\parskip}{6pt}
\title{QUBO Formulation: Maximum Independent Set Problem}
\date{}
\begin{document}
\maketitle

\subsection*{Maximum Independent Set Problem}

\paragraph{Problem Statement.}
Let $G=(V,E)$ be an undirected graph with vertex set $V$ and edge set $E$. An independent set is a subset $S \subseteq V$ such that no two vertices in $S$ are connected by an edge. The Maximum Independent Set Problem seeks to find an independent set $S$ of maximum cardinality, i.e., $\max |S|$ subject to the constraint that for every edge $(i,j) \in E$, at most one of $i$ or $j$ belongs to $S$.

\paragraph{Decision Variables.}
For each vertex $i \in V$, we introduce a binary decision variable $x_i \in \{0,1\}$. The variable $x_i$ takes the value $1$ if vertex $i$ is selected for the independent set, and $0$ otherwise. The decision vector is denoted $\mathbf{x} = (x_0, x_1, \dots, x_{n-1}) \in \{0,1\}^n$, where $n = |V|$.

\paragraph{QUBO Derivation.}
\textbf{Step 1 --- Objective Function.} The goal is to maximize the number of selected vertices, $\sum_{i \in V} x_i$. Since QUBO solvers minimize energy, we negate the objective:
\begin{align}
\min_{\mathbf{x} \in \{0,1\}^n} \left( -\sum_{i \in V} x_i \right).
\end{align}

\textbf{Step 2 --- Constraint Formulation.} For each edge $(i,j) \in E$, the independence condition requires that we cannot select both endpoints simultaneously. This is enforced by the linear inequality:
\begin{align}
x_i + x_j \le 1 \quad \forall (i,j) \in E.
\end{align}

\textbf{Step 3 --- Penalty Term Expansion.} Each inequality is converted to a quadratic penalty term by squaring the violation. For a given edge $(i,j) \in E$ and penalty weight $\lambda > 0$, the penalty is:
\begin{align}
P_{(i,j)}(\mathbf{x}) &= \lambda (x_i + x_j - 1)^2 \\
&= \lambda (x_i^2 + x_j^2 + 1 + 2x_i x_j - 2x_i - 2x_j).
\end{align}
Using the idempotent property of binary variables, $x_k^2 = x_k$ for all $k \in V$, the expression simplifies to:
\begin{align}
P_{(i,j)}(\mathbf{x}) &= \lambda (x_i + x_j + 1 + 2x_i x_j - 2x_i - 2x_j) \\
&= \lambda (1 - x_i - x_j + 2x_i x_j).
\end{align}
This is the general algebraic contribution of edge $(i,j)$ to the Hamiltonian.

\textbf{Step 4 --- Combined Hamiltonian.} Summing the negated objective and all edge penalties yields the total QUBO energy function $H(\mathbf{x})$:
\begin{align}
H(\mathbf{x}) &= -\sum_{i \in V} x_i + \sum_{(i,j) \in E} \lambda (1 - x_i - x_j + 2x_i x_j).
\end{align}

\textbf{Step 5 --- Q Matrix and Offset.} Collecting coefficients for linear and quadratic terms, and noting that each vertex $i$ appears in $\deg(i)$ edges, we obtain the closed-form QUBO parameters:
\begin{align}
Q_{i,i} &= -1 - \lambda \deg(i) \quad \forall i \in V, \\
Q_{i,j} &= 2\lambda \quad \forall (i,j) \in E, \\
Q_{i,j} &= 0 \quad \forall (i,j) \notin E, \\
c &= \lambda |E|.
\end{align}
The Hamiltonian is thus expressed as $H(\mathbf{x}) = \sum_{i \in V} Q_{i,i} x_i + \sum_{i < j} Q_{i,j} x_i x_j + c$.

\paragraph{Penalty Weight Justification.}
The penalty weight $\lambda$ must be chosen large enough to ensure that any violation of the independence constraint increases the total energy relative to the best feasible assignment. Violating a single edge constraint by setting $x_i = x_j = 1$ instead of the optimal valid configuration $(1,0)$ or $(0,1)$ yields a maximum objective gain of $1$ (selecting one extra vertex). The penalty incurred for this violation is $\lambda(1+1-1)^2 = \lambda$. Therefore, any $\lambda > 1$ strictly penalizes infeasible configurations more than the objective function rewards them. We select $\lambda = 2$ to provide a robust margin above the theoretical lower bound.

\paragraph{Correctness Argument.}
Minimizing $H(\mathbf{x})$ over $\{0,1\}^n$ recovers the optimal maximum independent set because (i) any feasible independent set satisfies $x_i + x_j \le 1$ for all $(i,j) \in E$, causing all penalty terms to vanish and reducing $H(\mathbf{x})$ to $-\sum_{i \in V} x_i$. Thus, among feasible solutions, minimizing $H$ is exactly equivalent to maximizing the set size. (ii) Any infeasible assignment violates at least one edge constraint, incurring a penalty of at least $\lambda$. Since $\lambda > 1$, the energy of any infeasible solution strictly exceeds that of the best feasible solution on the violated edge(s). Consequently, the global minimum of $H$ must be feasible and correspond to the maximum independent set.

\paragraph{Illustrative Example.}
Consider a path graph with $n=3$ vertices $V=\{0,1,2\}$ and edges $E=\{(0,1), (1,2)\}$. The degrees are $\deg(0)=1$, $\deg(1)=2$, and $\deg(2)=1$. Using $\lambda=2$, the general formulas from Step 5 yield the concrete QUBO parameters:
\begin{align}
Q_{0,0} &= -1 - 2(1) = -3, \\
Q_{1,1} &= -1 - 2(2) = -5, \\
Q_{2,2} &= -1 - 2(1) = -3, \\
Q_{0,1} &= 2(2) = 4, \\
Q_{1,2} &= 2(2) = 4, \\
c &= 2(2) = 4.
\end{align}
The resulting Hamiltonian is:
\begin{align}
H(\mathbf{x}) = 4 - 3x_0 - 5x_1 - 3x_2 + 4x_0 x_1 + 4x_1 x_2.
\end{align}
Evaluating $H$ over all $2^3=8$ binary assignments, the feasible independent sets are $(1,0,1)$, $(1,0,0)$, and $(0,0,1)$. Their energies are:
\begin{align}
H(1,0,1) &= 4 - 3 - 0 - 3 + 0 + 0 = -2, \\
H(1,0,0) &= 4 - 3 - 0 - 0 + 0 + 0 = 1, \\
H(0,0,1) &= 4 - 0 - 0 - 3 + 0 + 0 = 1.
\end{align}
Infeasible assignments such as $(1,1,0)$ incur a penalty of $2(1+1-1)^2 = 2$, yielding $H(1,1,0) = -2 + 2 = 0$. The global minimum is $H^* = -2$ at $\mathbf{x}^* = (1,0,1)$, correctly identifying the maximum independent set $\{0,2\}$ of size $2$.

\end{document}